% ============================================================
%  SmartRide : Project Phase 1 Report
%  B.Tech CSE, TKM Institute of Technology (APJ Abdul Kalam
%  Technological University)
%  Group 14 -- Interim / Phase 1 Report
% ============================================================
\documentclass[12pt,a4paper,oneside]{report}

\usepackage[left=1.25in,right=1.25in,top=1in,bottom=1in]{geometry}
\usepackage[T1]{fontenc}
\usepackage[utf8]{inputenc}
\usepackage{setspace}
\usepackage{graphicx}
\usepackage{amsmath,amssymb}
\usepackage{array}
\usepackage{longtable}
\usepackage{multirow}
\usepackage{enumitem}
\usepackage{titlesec}
\usepackage{tocloft}
\usepackage{tikz}
\usepackage{pgfplots}
\usepackage[hidelinks,plainpages=false,pdfpagelabels]{hyperref}
\usepackage{caption}
\usepackage{float}
\newcommand{\rupee}{\textrm{Rs.}\,}
\newcommand{\SI}[2]{#1\,#2}

\usetikzlibrary{shapes.geometric,arrows.meta,positioning,calc,fit,backgrounds,decorations.pathreplacing}
\pgfplotsset{compat=1.17}

\onehalfspacing
\setlength{\emergencystretch}{2em}
\hyphenation{multi-con-stel-la-tion ac-cel-er-om-e-ter mi-cro-con-trol-ler}
\setlength{\parindent}{1.5em}
\setlength{\parskip}{0pt}

% ---------- Chapter / Section heading style ----------
\titleformat{\chapter}[display]
  {\centering\normalfont\large\bfseries}
  {CHAPTER \thechapter}{1.2em}{\large}
\titlespacing*{\chapter}{0pt}{-20pt}{35pt}

\titleformat{\section}
  {\normalfont\large\bfseries}{\thesection}{0.8em}{}
\titlespacing*{\section}{0pt}{2.2ex plus 1ex minus .2ex}{1.6ex plus .2ex}

\titleformat{\subsection}
  {\normalfont\normalsize\bfseries}{\thesubsection}{0.8em}{}
\titlespacing*{\subsection}{0pt}{2ex plus 1ex minus .2ex}{1.3ex plus .2ex}

% ---------- Table of contents style ----------
\renewcommand{\cftchapfont}{\bfseries}
\renewcommand{\cftsecfont}{\bfseries}
\renewcommand{\cftsubsecfont}{\bfseries}
\renewcommand{\cftchappagefont}{\bfseries}
\setlength{\cftbeforechapskip}{0.6em}

% ---------- Unnumbered centred front-matter heading ----------
\newcommand{\frontheading}[1]{%
  \begin{center}\large\bfseries\MakeUppercase{#1}\end{center}
  \vspace{2.2em}}

% ---------- Handy column types ----------
\newcolumntype{L}[1]{>{\raggedright\arraybackslash}p{#1}}
\newcolumntype{J}[1]{>{\arraybackslash}p{#1}}

% ---------- Verdict / bullet helpers ----------
\newcommand{\verdict}[1]{\par\noindent\textbf{Verdict:} #1\par}

% ---------- Always prefix chapter number (gives Table 0.1 in front matter) ----------
\renewcommand{\thetable}{\thechapter.\arabic{table}}
\renewcommand{\thefigure}{\thechapter.\arabic{figure}}

\begin{document}
\pagenumbering{Roman}
\pagestyle{empty}

% ============================================================
%  TITLE PAGE
% ============================================================
\thispagestyle{empty}
\begin{singlespace}
\begin{center}
\vspace*{0.45in}
{\large\bfseries SMARTRIDE: BLUETOOTH-BASED TWO-WHEELER\\[0.45em]
RIDE LOGGER AND CRASH DETECTION SYSTEM\\[0.45em]
USING ESP32-S3\par}

\vspace{0.45in}
{\large A PROJECT REPORT\par}
\vspace{0.14in}
{submitted by\par}

\vspace{0.3in}
{\renewcommand{\arraystretch}{1.35}
\begin{tabular}{@{}p{2.8in}@{}p{1.6in}@{}}
\textbf{Adarsh Raj}        & (TKI23CS005)\\
\textbf{Muhammed Afsal}    & (TKI23CS074)\\
\textbf{Irfan V}           & (TKI23CS050)\\
\textbf{Rojith Ram}        & (TKI23CS097)\\
\end{tabular}}

\vspace{0.28in}
{to\par}
\vspace{0.16in}
{the \textbf{APJ Abdul Kalam Technological University}\par}
\vspace{0.14in}
{in partial fulfillment of the requirements for the award of the Degree of\par}
\vspace{0.16in}
{\textbf{Bachelor of Technology} in\par}
\vspace{0.12in}
{\textit{Computer Science and Engineering}\par}

\vspace{0.22in}
\includegraphics[width=1.5in]{tkmlogo.png}

\vspace{0.3in}
{\textbf{Department of Computer Science and Engineering}\par}
\vspace{0.18in}
{\textbf{TKM Institute of Technology}\par}
\vspace{0.16in}
{Karuvelil, Kollam\par}

\vspace{0.35in}
{September 2026\par}
\end{center}
\end{singlespace}
\clearpage

% ============================================================
%  DECLARATION
% ============================================================
\thispagestyle{empty}
\frontheading{Declaration}

We undersigned hereby declare that the project report ``SmartRide: Bluetooth-Based
Two-Wheeler Ride Logger and Crash Detection System Using ESP32-S3'', submitted for
partial fulfillment of the requirements for the award of degree of Bachelor of Technology
of the APJ Abdul Kalam Technological University, Kerala, is a bonafide work done by us
under the supervision of Mrs. Neetha Alex.

This submission represents our ideas in our own words and where ideas or words of others
have been included, we have adequately and accurately cited and referenced the original
sources. We also declare that we have adhered to ethics of academic honesty and integrity
and have not misrepresented or fabricated any data or idea or fact or source in our
submission. We understand that any violation of the above will be a cause for disciplinary
action by the Institute and/or the University and can also evoke penal action from the
sources which have thus not been properly cited or from whom proper permission has not
been obtained.

This report has not been previously formed the basis for the award of any degree, diploma
or similar title of any other University.

\vspace{1.1in}
\noindent
\begin{minipage}[t]{0.45\textwidth}
Place: Kollam\\[1.2em]
Date:
\end{minipage}%
\begin{minipage}[t]{0.5\textwidth}
Adarsh Raj\\
Muhammed Afsal\\
Irfan V\\
Rojith Ram
\end{minipage}
\clearpage

% ============================================================
%  CERTIFICATE
% ============================================================
\thispagestyle{empty}
\begin{center}
{\normalsize\bfseries DEPARTMENT OF COMPUTER SCIENCE AND ENGINEERING\par}
\vspace{0.9em}
{\normalsize\bfseries TKM INSTITUTE OF TECHNOLOGY, KARUVELIL, KOLLAM\par}

\vspace{0.28in}
\includegraphics[width=1.3in]{tkmlogo.png}

\vspace{0.28in}
{\large\bfseries CERTIFICATE\par}
\end{center}
\vspace{1.4em}

This is to certify that the report entitled ``\textbf{SmartRide: Bluetooth-Based
Two-Wheeler Ride Logger and Crash Detection System Using ESP32-S3}'' submitted by
\textbf{Adarsh Raj}, \textbf{Muhammed Afsal}, \textbf{Irfan V}, \textbf{Rojith Ram} to the
APJ Abdul Kalam Technological University in partial fulfilment of the requirements for the
award of degree of Bachelor of Technology in Computer Science and Engineering is a Bonafide
record of the project carried out by them under our guidance and supervision. This report in
any form has not been submitted to any other University or Institute for any purpose.

\vspace{1.0in}
\begin{singlespace}
\renewcommand{\arraystretch}{1.25}
\setlength{\tabcolsep}{0pt}
\noindent
\begin{tabular}{@{\hspace{0.45in}}p{0.52\textwidth}@{}p{0.40\textwidth}@{}}
\textbf{Mrs. Neetha Alex} & \textbf{Dr. Nijil Raj N}\\
Assistant Professor       & Head of the Department\\
Project Guide             & \\[0.85in]
\textbf{Dr. Reji R}       & \textbf{Mr. Rupesh Ravi M R}\\
Professor                 & Assistant Professor\\
Project Coordinator       & Project Coordinator\\
\end{tabular}
\end{singlespace}
\clearpage

% ============================================================
%  ACKNOWLEDGEMENT
% ============================================================
\thispagestyle{empty}
\frontheading{Acknowledgement}

First of all, We thank God Almighty for helping us to successfully complete our Project.

We express our sincere thanks and a deep sense of humble gratitude to the Principal of this
institution, \textbf{Dr. Gouri Mohan L}, for providing all the necessary facilities.

We thank \textbf{Dr. Nijil Raj N}, Head of the Department, Computer Science \& Engineering,
for all the words of inspiration.

Next, We would like to thank our project coordinators \textbf{Dr. Reji R},
\textbf{Mr. Rupesh Ravi} for their valuable advices.

We have this great pleasure to express our deep sense of gratitude and obligation to our
project guide \textbf{Mrs. Neetha Alex} for her valuable guidance and suggestions
throughout the entire project work.

Last but certainly not the least, We would also like to thank all faculty and staff of the
Department of Computer Science \& Engineering and my friends for their help and
co-operation.

\vspace{1.0in}
\begin{singlespace}
\renewcommand{\arraystretch}{2.0}
\noindent\hspace*{0.40\textwidth}%
\begin{tabular}{@{}l@{\hspace{1.2em}}l@{}}
\textbf{TKI23CS005} & \textbf{Adarsh Raj}\\
\textbf{TKI23CS074} & \textbf{Muhammed Afsal}\\
\textbf{TKI23CS050} & \textbf{Irfan V}\\
\textbf{TKI23CS097} & \textbf{Rojith Ram}
\end{tabular}
\end{singlespace}
\clearpage

% ============================================================
%  ABSTRACT
% ============================================================
\thispagestyle{empty}
\frontheading{Abstract}

Two-wheelers constitute the largest share of registered vehicles on Indian roads, yet the
overwhelming majority of them offer the rider no digital record of how they are ridden and
no assistance when the rider falls. This project, \textit{SmartRide: Bluetooth-Based
Two-Wheeler Ride Logger and Crash Detection System Using ESP32-S3}, proposes a compact,
self-contained on-board unit that logs ride data, recognises crashes from motion signatures
and delivers both to the rider's smartphone over Bluetooth Low Energy, without depending on
cellular connectivity or cloud services.

The on-board unit is built around an ESP32-S3 microcontroller, which reads location, speed
and course from an ATGM336H-5N multi-constellation GNSS receiver over UART and acquires
motion data from an MPU6500 six-axis accelerometer and gyroscope over a high-speed serial
interface. An ADXL345 three-axis accelerometer is provided as an optional secondary impact
channel, offering an independent hardware interrupt for corroborating high-energy events.
The unit is powered from a lithium-ion cell managed by a dedicated battery management
system (BMS), which enforces over-charge, over-discharge and over-current protection and
allows the logger to remain alive through the shock and supply interruption of a crash.

Ride parameters such as instantaneous speed, cumulative distance, duration and travelled
route are computed on the microcontroller itself and exposed as Bluetooth Low Energy GATT
characteristics. Crash detection uses a multi-condition rule that combines resultant
acceleration, sustained tilt derived from gyroscope integration, and post-impact stillness,
so that potholes, speed breakers and kerb strikes do not raise false alarms. A Flutter
Android application acts purely as a BLE client: it renders a live dashboard, plots the
route, stores ride history in a local SQLite database and issues distance-based maintenance
reminders from the cumulative odometer reading.

By keeping sensing, computation, crash logic and storage entirely on the device and the
paired handset, SmartRide delivers an affordable, network-independent and privacy-preserving
smart riding solution. The system requires no SIM card, no recurring data plan and no
vendor cloud account, which makes it suitable as a low-cost retrofit for the existing
two-wheeler population.

\vspace{1.2em}
\noindent\textbf{Keywords:}~\textit{Internet of Things, Two-Wheeler Telematics, Ride
Logging, Crash Detection, Bluetooth Low Energy, ESP32-S3, MPU6500, ATGM336H-5N, ADXL345,
GNSS, Flutter, Offline Operation, Rider Safety}
\clearpage

% ============================================================
%  MAPPING WITH PROGRAM OUTCOMES
% ============================================================
\thispagestyle{empty}
\begin{center}{\large\bfseries MAPPING WITH PROGRAM OUTCOMES (POs)\par}\end{center}
\vspace{2.2em}

The mapping between the project contributions and relevant POs is summarized below.

\vspace{1.2em}
\begin{center}
\begin{singlespace}
\small
\renewcommand{\arraystretch}{1.3}
\begin{tabular}{|L{1.3in}|J{3.5in}|}
\hline
\textbf{PO} & \textbf{Contribution of the Project}\\
\hline
PO1: Engineering Knowledge & Applied principles of embedded systems, sensor signal
processing, wireless communication and mobile application development to realise a
self-contained two-wheeler ride logger.\\
\hline
PO2: Problem Analysis & Identified and analysed the absence of affordable ride telematics
and crash notification on two-wheelers, and the dependence of existing solutions on
cellular coverage.\\
\hline
PO3: Design/Development of Solutions & Designed a complete on-board architecture comprising
the ESP32-S3, GNSS receiver, inertial sensors, battery management and a BLE-connected
Android application.\\
\hline
PO4: Conduct Investigations of Complex Problems & Planned controlled drop, tilt and on-road
trials to characterise sensor behaviour and to calibrate crash-detection thresholds against
false-alarm sources.\\
\hline
PO5: Modern Tool Usage & Utilised ESP-IDF, the Arduino core for ESP32-S3, Flutter, SQLite,
Python-based log analysis and Git for efficient development, validation and version
control.\\
\hline
PO6: The Engineer and Society & Addresses a road-safety problem of direct public
consequence by shortening the interval between a two-wheeler crash and the moment help is
summoned.\\
\hline
PO7: Environment and Sustainability & Encourages timely, distance-based servicing that
extends vehicle life and reduces emissions from poorly maintained engines, supporting
SDG 11 and SDG 12.\\
\hline
PO9: Individual and Team Work & Enabled the team to divide firmware, application, hardware
integration and testing responsibilities while working towards a single integrated
deliverable.\\
\hline
PO10: Communication & Developed the ability to document and present technical work through
interim evaluation, design reviews and this project report.\\
\hline
PO12: Life-long Learning & Required independent study of GNSS protocols, inertial sensing,
the Bluetooth Low Energy GATT model and lithium-ion battery management beyond the
prescribed syllabus.\\
\hline
\end{tabular}
\captionof{table}{Mapping of Project with Program Outcomes (POs)}
\end{singlespace}
\end{center}
\clearpage

% ============================================================
%  MAPPING WITH COURSE OUTCOMES
% ============================================================
\thispagestyle{empty}
\begin{center}{\large\bfseries MAPPING WITH COURSE OUTCOMES (COs)\par}\end{center}
\vspace{2.2em}

The mapping between the project objectives and relevant Course Outcomes (COs) is
summarized below.

\vspace{1.4em}
\begin{center}
\begin{singlespace}
\small
\renewcommand{\arraystretch}{1.3}
\begin{tabular}{|L{1.3in}|J{3.5in}|}
\hline
\textbf{CO} & \textbf{Contribution of the Project}\\
\hline
CO1: Problem Identification & Identified real-world problems in smart transportation --- the
absence of ride history, delayed crash response and guesswork-based servicing --- and
proposed an IoT-based solution. \emph{(Mapped to PO1, PO2)}\\
\hline
CO2: System Design & Designed and specified an ESP32-S3 based Bluetooth ride-logging system
with defined hardware interfaces, firmware tasks and a BLE GATT service.
\emph{(Mapped to PO3)}\\
\hline
CO3: Implementation & Integrated the ATGM336H-5N GNSS receiver, the MPU6500 inertial
measurement unit and BLE communication for real-time ride monitoring.
\emph{(Mapped to PO5)}\\
\hline
CO4: Data Analysis & Analysed logged ride and impact data to evaluate distance accuracy,
detection latency and false-alarm behaviour through field testing.
\emph{(Mapped to PO4)}\\
\hline
CO5: Use of Modern Tools & Applied ESP-IDF, Flutter, SQLite, the Google Maps API and Python
analysis scripts for development, visualisation and threshold tuning.
\emph{(Mapped to PO5)}\\
\hline
CO6: Societal Benefit and Sustainability & Delivered a safety-oriented, affordable solution
that benefits riders and promotes timely maintenance.
\emph{(Mapped to PO6, PO7)}\\
\hline
CO7: Teamwork and Communication & Worked collaboratively with clear task allocation and
documented the project professionally through reviews and reports.
\emph{(Mapped to PO9, PO10, PO12)}\\
\hline
\end{tabular}
\captionof{table}{Mapping of Project with Course Outcomes (COs)}
\end{singlespace}
\end{center}
\clearpage

% ============================================================
%  MAPPING WITH SDGs
% ============================================================
\thispagestyle{empty}
\begin{center}{\large\bfseries MAPPING WITH SUSTAINABLE DEVELOPMENT\\[0.25em]
GOALS (SDGs)\par}\end{center}
\vspace{2.2em}

The project \textbf{SmartRide: Bluetooth-Based Two-Wheeler Ride Logger and Crash Detection
System Using ESP32-S3} has strong alignment with the United Nations Sustainable Development
Goals (SDGs). The mapping of the project to the relevant SDGs is presented below.

\vspace{1.6em}
\begin{center}
\begin{singlespace}
\small
\renewcommand{\arraystretch}{1.3}
\begin{tabular}{|L{1.3in}|J{3.5in}|}
\hline
\textbf{SDG} & \textbf{Contribution of the Project}\\
\hline
SDG 3: Good Health and Well-being & Target 3.6 aims to reduce deaths and injuries from road
traffic accidents. Faster crash notification shortens the interval before help reaches an
injured two-wheeler rider.\\
\hline
SDG 9: Industry, Innovation and Infrastructure & Demonstrates affordable embedded innovation
that can be retrofitted to vehicles already on the road, requiring no new roadside or
network infrastructure.\\
\hline
SDG 11: Sustainable Cities and Communities & Target 11.2 seeks safe and accessible transport
for all. Ride data and automatic safety alerts contribute directly to safer urban
mobility.\\
\hline
SDG 12: Responsible Consumption and Production & Distance-based maintenance reminders
promote timely servicing, extending vehicle life and reducing emissions and waste from
poorly maintained engines.\\
\hline
\end{tabular}
\captionof{table}{Mapping of Project with Sustainable Development Goals (SDGs)}
\end{singlespace}
\end{center}
\clearpage

% ============================================================
%  FEASIBILITY REPORT
% ============================================================
\thispagestyle{empty}
\begin{center}{\large\bfseries FEASIBILITY REPORT\par}\end{center}
\vspace{2.0em}

\noindent\hspace{0.6em}{\large\bfseries (i) TECHNICAL FEASIBILITY}\par
\vspace{0.4em}
\noindent\hspace{0.6em}\textbf{Hardware Feasibility:}
\begin{itemize}[leftmargin=*,itemsep=0.35em,topsep=0.5em]
\item \textbf{Microcontroller:} The ESP32-S3 provides a dual-core Xtensa LX7 processor at up
to \SI{240}{MHz}, ample SRAM with optional PSRAM, and an integrated Bluetooth~5 Low Energy
radio --- sufficient for concurrent sensor sampling, on-device computation and BLE service
hosting.

\item \textbf{Positioning:} The ATGM336H-5N is a low-cost multi-constellation GNSS module
that outputs standard NMEA sentences over UART and provides a 1\,PPS timing output, giving
reliable position, speed and course fixes under Indian sky conditions.

\item \textbf{Motion Sensing:} The MPU6500 six-axis accelerometer and gyroscope supports
selectable ranges up to $\pm 16\,g$ and $\pm 2000^{\circ}/\mathrm{s}$, an on-chip
FIFO and a wake-on-motion interrupt, which together make impact and tilt detection
practical at low power.

\item \textbf{Optional Redundant Impact Channel:} The ADXL345 three-axis accelerometer
offers built-in tap, activity and free-fall interrupt engines on an independent bus address,
allowing a second opinion on high-energy events at negligible cost.

\item \textbf{Power and Protection:} A lithium-ion cell with a dedicated battery management
system (BMS) --- already available with the team --- supplies the unit and enforces
over-charge, over-discharge, over-current and short-circuit protection, so the logger
survives supply interruption during a crash.
\end{itemize}
\verdict{All components are inexpensive, readily available in the Indian market,
electrically compatible at \SI{3.3}{V} logic, and require only moderate integration effort.}

\vspace{0.8em}
\noindent\hspace{0.6em}\textbf{Software Feasibility:}
\begin{itemize}[leftmargin=*,itemsep=0.35em,topsep=0.5em]
\item \textbf{Development Tools:} ESP-IDF and the Arduino core for the ESP32-S3, both
mature and fully documented for this part.
\item \textbf{Mobile Application:} Flutter for Android, with well-maintained BLE and
mapping plugins.
\item \textbf{On-device Storage:} SQLite within the mobile application for ride history,
and non-volatile flash on the microcontroller for buffered records.
\item \textbf{Analysis:} Python with NumPy and Matplotlib for offline log analysis and
threshold tuning.
\end{itemize}
\verdict{Every tool in the chain is open-source or free for academic use, keeping
development cost effectively zero.}

\vspace{0.6em}
\noindent\hspace{0.6em}\textbf{System Integration:} Sensor data flows from the GNSS receiver
and the inertial measurement unit into the ESP32-S3, where it is filtered, fused into ride
statistics and evaluated against the crash rule. Results are published as BLE GATT
characteristics and consumed by the Flutter application, which persists them to SQLite and
raises a cancellable crash alert. No component of the critical path leaves the vehicle or
the paired handset.
\verdict{Integration is technically achievable with the team's existing skill set, and the
modular design leaves room for later extension.}
\clearpage

% ============================================================
%  ECONOMIC FEASIBILITY
% ============================================================
\thispagestyle{empty}
\noindent\hspace{0.6em}{\large\bfseries (ii) ECONOMIC FEASIBILITY}\par
\vspace{1.1em}

\noindent\textbf{Prototype Development Cost:}
\begin{itemize}[leftmargin=*,itemsep=0.3em,topsep=0.5em]
\item ESP32-S3 development board: \(\rupee\)450--700
\item ATGM336H-5N GNSS module with active antenna: \(\rupee\)250--400
\item MPU6500 inertial measurement unit module: \(\rupee\)120--220
\item ADXL345 accelerometer module (optional secondary channel): \(\rupee\)120--200
\item Lithium-ion cell with BMS: \emph{already available with the team} (indicative market
value \(\rupee\)300--500)
\item Enclosure, mounting hardware, wiring harness and connectors: \(\rupee\)250--450
\item Miscellaneous (perfboard or PCB, headers, passives, fasteners): \(\rupee\)200--350
\end{itemize}

\noindent\textbf{Total Incremental Hardware Cost:} \(\rupee\)1{,}270--2{,}120 per unit
(excluding the battery and BMS already in hand).\\
\noindent\textbf{Total Build Cost if Purchased Outright:} \(\rupee\)1{,}570--2{,}620 per
unit.

\vspace{1.0em}
\noindent\textbf{Software Cost:} Nil. ESP-IDF, the Arduino core, Flutter, SQLite and Python
are all free; the Google Maps API remains within its free usage tier at the scale of a
single prototype.
\begin{itemize}[leftmargin=*,itemsep=0.3em,topsep=0.5em]
\item No SIM card, no cellular data plan and no recurring subscription of any kind.
\item No cloud hosting or database charges, since all storage is on the device and the
paired handset.
\end{itemize}

\vspace{0.8em}
\noindent\textbf{Cost-Benefit Perspective:} Commercial two-wheeler telematics devices in
India typically cost several thousand rupees and additionally require a monthly data plan.
SmartRide removes the recurring cost entirely and keeps the one-time cost within the price
of a routine service. Against this, the rider gains a permanent ride record, automatic
crash alerting and maintenance reminders tied to actual distance travelled. The low unit
cost and the absence of any subscription make the design attractive both as an academic
prototype and as a candidate for mass-market retrofit.
\clearpage

% ============================================================
%  FRONT MATTER LISTS
% ============================================================
\pagestyle{plain}
\pagenumbering{roman}
\setcounter{page}{1}

\tableofcontents
\clearpage

\listoftables
\clearpage

\listoffigures
\clearpage

% ============================================================
%  ABBREVIATIONS
% ============================================================
\begin{center}{\large\bfseries ABBREVIATIONS\par}\end{center}
\vspace{1.6em}
\begin{singlespace}
\noindent
\begin{tabular}{@{}p{1.15in}@{}p{4.0in}@{}}
ADC       & Analog-to-Digital Converter\\
ADXL345   & Analog Devices Three-Axis Digital Accelerometer, Model 345\\
API       & Application Programming Interface\\
ATGM336H  & Zhongke Micro Multi-Constellation GNSS Module Series\\
BDS       & BeiDou Navigation Satellite System\\
BLE       & Bluetooth Low Energy\\
BMS       & Battery Management System\\
CEP       & Circular Error Probable\\
DFD       & Data Flow Diagram\\
DMP       & Digital Motion Processor\\
ESP32-S3  & Espressif 32-bit Xtensa LX7 Dual-Core SoC with Wi-Fi and BLE\\
ESP-IDF   & Espressif IoT Development Framework\\
FIFO      & First In, First Out (sensor sample buffer)\\
GATT      & Generic Attribute Profile\\
GLONASS   & Global Navigation Satellite System (Russia)\\
GNSS      & Global Navigation Satellite System\\
GPIO      & General Purpose Input/Output\\
GPS       & Global Positioning System\\
GSM       & Global System for Mobile Communications\\
IDE       & Integrated Development Environment\\
IMU       & Inertial Measurement Unit\\
IoT       & Internet of Things\\
I\textsuperscript{2}C & Inter-Integrated Circuit (serial bus)\\
MPU6500   & InvenSense Six-Axis Motion Tracking Device, Model 6500\\
NavIC     & Navigation with Indian Constellation\\
NMEA      & National Marine Electronics Association (sentence format)\\
NVS       & Non-Volatile Storage\\
PPS       & Pulse Per Second\\
PSRAM     & Pseudo-Static Random Access Memory\\
RTOS      & Real-Time Operating System\\
SDG       & Sustainable Development Goal\\
SoC       & System on Chip\\
SPI       & Serial Peripheral Interface\\
SQLite    & Embedded Relational Database Engine\\
UART      & Universal Asynchronous Receiver/Transmitter\\
UUID      & Universally Unique Identifier\\
\end{tabular}
\end{singlespace}
\clearpage

\pagenumbering{arabic}
\setcounter{page}{1}
% ============================================================
\chapter{INTRODUCTION}

\section{BACKGROUND AND CONTEXT}

Two-wheelers are the dominant mode of personal transport in India. They account for the
largest share of registered vehicles in the country, and in a state such as Kerala they are
the everyday vehicle of students, daily commuters, delivery riders and small traders alike.
Their affordability, fuel economy and ability to negotiate congested roads make them
indispensable. The same characteristics, however, leave the rider far more exposed than the
occupant of a car: there is no crumple zone, no seat belt and no airbag, and a fall at even
moderate speed can be serious.

Despite this exposure, the two-wheeler has remained almost entirely undigitised. A modern
car routinely records trip data, warns about servicing, and in many markets places an
automatic emergency call after a collision. A typical motorcycle or scooter offers none of
this. The odometer shows a single cumulative number, the trip meter is reset by hand, and
the rider has no record of where a ride went, how fast it was, how long it took or how
harshly the vehicle was handled. Servicing intervals are recalled from memory or estimated
from the last workshop visit rather than derived from distance actually travelled.

The safety gap is more consequential still. When a rider falls on a quiet stretch of road
--- and a large share of two-wheeler crashes in India occur away from dense traffic --- there
may be nobody present to raise an alarm. The interval between the impact and the moment
somebody is informed is the single most important determinant of the outcome of a serious
injury, and it is precisely this interval that no low-cost consumer product currently
addresses for two-wheelers.

Commercial telematics solutions do exist, but their design assumptions make them a poor fit
for this problem. They are built around a cellular modem, which brings with it the cost of
the module itself, a SIM card, a recurring data plan and dependence on network coverage that
is unreliable on exactly the rural roads where help is most needed. They also route personal
movement data through a vendor cloud, which raises both privacy questions and the prospect
of the device becoming useless if the service is discontinued. The result is a category of
product priced and provisioned for fleet operators rather than for the individual rider.

At the same time, the components required to solve the problem locally have become
inexpensive and capable. Modern microcontrollers such as the Espressif ESP32-S3 combine a
dual-core processor, generous memory and a Bluetooth~5 Low Energy radio in a single
low-cost module. Multi-constellation satellite receivers such as the ATGM336H-5N deliver
usable position and speed fixes for a few hundred rupees. Six-axis inertial measurement
units such as the MPU6500 provide acceleration and angular-rate data at rates and ranges
that are entirely adequate for recognising an impact and a fall. Lithium-ion cells with
integrated battery management have become a standard, safe and readily available building
block. Meanwhile, every rider already carries a capable computer with a colour display, a
map application and persistent storage --- a smartphone.

\textit{SmartRide} is built on the observation that these facts, taken together, remove the
need for the cellular link altogether. If the on-board unit performs its own sensing,
filtering, ride computation and crash reasoning, and if the smartphone in the rider's pocket
serves as the display and the archive, then the only communication link required is a short
Bluetooth Low Energy connection between the two. There is no SIM card, no data plan, no
cloud account and no dependence on network coverage. The system works as well on an
unlit rural road as it does in the middle of a city, and the rider's movement history never
leaves the two devices they own.

\section{PROBLEM STATEMENT}

Riders of two-wheelers in India have no affordable, self-contained means of recording how
their vehicle is ridden and of summoning help automatically when they fall.

The absence of a ride record deprives the rider of speed, distance, duration and route
history for any journey, and therefore of any objective basis for reviewing their own riding
or substantiating it to an insurer or employer. The absence of automatic crash notification
means that a fall on an isolated road can remain unnoticed until somebody happens to pass,
so that help is delayed at precisely the point where delay matters most. The absence of
distance-linked maintenance tracking means that servicing is scheduled by guesswork rather
than by actual kilometres covered, which shortens the life of the vehicle and increases both
its running cost and its emissions.

Existing telematics products do not close this gap for the individual rider. They are
expensive, they require a SIM card and a recurring data plan, they depend on cellular
coverage that fails on rural roads, and they lock the user's movement data into a
proprietary cloud service. Conventional academic accident-alert prototypes based on
GSM or GPRS inherit the same connectivity dependence, and additionally tend to trigger on
any large acceleration, producing false alarms from potholes, speed breakers and kerb
strikes that quickly destroy the rider's trust in the device.

There is therefore a need for a low-cost, self-contained on-board unit that logs ride data,
recognises genuine crashes from motion signatures with an acceptably low false-alarm rate,
and delivers this information to the rider's smartphone over Bluetooth Low Energy, without
relying on cellular connectivity or subscription services.

\section{OBJECTIVES}

The objectives of this project are as follows:

\begin{itemize}[leftmargin=*,itemsep=0.7em,topsep=0.7em]
\item \textbf{To design a compact ESP32-S3 based on-board unit powered from a lithium-ion
cell with an integrated battery management system.} The ESP32-S3 was selected for its
dual-core Xtensa LX7 processor, its generous on-chip memory and its native Bluetooth~5 Low
Energy radio, which together allow sensor acquisition, ride computation and BLE service
hosting to run concurrently on one inexpensive module. Powering the unit from its own
managed cell rather than directly from the vehicle supply keeps the logger alive through the
electrical disturbance of a crash.

\item \textbf{To acquire location, speed, course and distance from an ATGM336H-5N
multi-constellation GNSS receiver.} The module delivers standard NMEA sentences over UART
together with a one-pulse-per-second timing output, from which the firmware derives
instantaneous speed, cumulative distance by the Haversine formula, and the travelled route
as a sequence of timestamped coordinates.

\item \textbf{To detect impacts and falls from MPU6500 accelerometer and gyroscope data.}
The six-axis inertial measurement unit is sampled at a fixed high rate through its on-chip
FIFO, and the firmware evaluates resultant acceleration together with orientation derived
from the gyroscope in order to distinguish a genuine fall from ordinary road disturbance.

\item \textbf{To provide an optional secondary impact channel using an ADXL345
accelerometer.} Where additional confidence is required, the ADXL345 can be populated on the
same bus at a separate address and configured to raise a hardware interrupt through its
built-in tap and activity engines, giving an independent corroboration of a high-energy
event without occupying processor time.

\item \textbf{To implement multi-condition crash detection that suppresses false alarms.}
Rather than triggering on acceleration alone, the firmware requires an impact, a sustained
tilt and a period of post-impact stillness to occur in sequence before a crash is declared,
so that potholes, speed breakers and kerb strikes are rejected.

\item \textbf{To stream ride data and crash events to an Android application over Bluetooth
Low Energy.} The on-board unit acts as a BLE GATT server exposing ride statistics, route
points and crash notifications as characteristics, and continues to log and buffer data
locally when the phone is out of range.

\item \textbf{To store ride history on the handset and generate distance-based maintenance
reminders.} The Flutter application persists each completed ride in a local SQLite database
and maintains a cumulative odometer from which servicing reminders are raised at
user-configurable intervals.

\item \textbf{To keep the entire system independent of cellular networks and cloud
services.} Every element of the critical path --- sensing, computation, crash reasoning,
alerting and storage --- resides either on the on-board unit or on the paired handset, so
that the system requires no SIM card, no data plan and no vendor account.

\item \textbf{To validate accuracy, latency and reliability through controlled and on-road
trials.} Distance and speed are to be compared against GNSS ground truth, crash detection
is to be characterised for detection latency and false-alarm rate, and the BLE link is to be
verified over typical rider-to-phone separation.
\end{itemize}

\section{SCOPE AND LIMITATIONS}

\subsection{Scope}

\begin{itemize}[leftmargin=*,itemsep=0.45em,topsep=0.6em]
\item The system targets a single two-wheeler paired with a single Android smartphone.

\item Fully offline operation: ride data and crash events are stored on the on-board unit
and on the handset, with no cloud or cellular component in the critical path.

\item Ride parameters covered are instantaneous speed, cumulative distance, ride duration,
travelled route and the timestamp of each record.

\item Crash detection from inertial data, with the resulting alert raised inside the mobile
application and cancellable by the rider within a countdown window.

\item An Android application developed with Flutter, comprising a live dashboard, a route
map, a searchable ride history backed by SQLite, and distance-based maintenance reminders.

\item Optional population of the ADXL345 as a redundant impact channel, selectable at build
time without any change to the application.
\end{itemize}

\subsection{Limitations}

\begin{itemize}[leftmargin=*,itemsep=0.45em,topsep=0.6em]
\item The unit does not place an emergency call or send an SMS to responders, because no
cellular modem is fitted; the alert is delivered to the paired handset only. Onward
notification to an emergency contact is identified as future work once a companion cellular
variant is considered.

\item Crash detection is threshold- and rule-based rather than learned. No machine-learning
classifier is trained in the current phase, and the thresholds are calibrated empirically
rather than derived from a large labelled crash dataset.

\item Because the design is deliberately offline, a crash cannot be reported if the paired
handset is absent, powered off or out of Bluetooth range at the moment of the event; the
unit will, however, retain the event in its local buffer and deliver it on the next
connection.

\item Positioning accuracy is limited by the GNSS module and by sky visibility. Performance
degrades in dense urban canyons, tunnels and heavy tree cover, and no dead-reckoning
fallback is implemented in this phase.

\item Only Android is supported. An iOS build, while feasible with the same Flutter
codebase, is outside the present scope.

\item The system does not interface with the vehicle's electronic control unit, braking
system or any factory electronics; it is a purely additive, non-intrusive retrofit.

\item Validation is limited to the number of controlled drop, tilt and on-road trials that
can be conducted safely by the team; no instrumented crash testing is performed.
\end{itemize}

\section{INDUSTRY RELEVANCE}

Connected-vehicle technology is one of the fastest-growing segments of the automotive
industry, and two-wheelers are its largest unserved market. Telematics has matured rapidly
in four-wheelers, where trip logging, usage-based insurance, predictive maintenance and
automatic collision notification are becoming standard. Very little of that capability has
reached the two-wheeler, largely because the cost structure of conventional telematics ---
cellular module, SIM, data plan and cloud back end --- cannot be supported at motorcycle
price points.

SmartRide is relevant to this transformation precisely because it inverts that cost
structure. By performing sensing, computation and crash reasoning at the edge and using the
rider's own smartphone as the display and archive, it demonstrates that the core value of
telematics --- a trustworthy ride record and automatic crash notification --- can be
delivered without any recurring cost. This edge-first, phone-as-gateway pattern is the same
one now being adopted across wearables, asset trackers and industrial condition-monitoring
products, and the architecture presented here is directly transferable to those domains.

Several industry segments map onto the work. The insurance sector is moving towards
usage-based and behaviour-based premiums, for which an auditable record of distance, speed
and harsh-event history is exactly the required input. Two-wheeler manufacturers and
accessory suppliers in India are actively adding connected features to attract younger
buyers, and a BLE-only module represents the least expensive route to that feature set.
Fleet operators in the rapidly expanding food-delivery and last-mile logistics sector run
large two-wheeler fleets on which crash detection and distance-accurate servicing have
immediate operational value. Vehicle service networks benefit from maintenance scheduling
driven by real distance rather than by calendar approximation.

The specific component choices also reflect current industry practice. The ESP32-S3 is
widely deployed in commercial connected products; multi-constellation receivers of the
ATGM336H class are standard in cost-sensitive tracking hardware; MEMS inertial measurement
units of the MPU6500 family are used across consumer and automotive-adjacent applications;
and lithium-ion packs with integrated battery management are the accepted approach to safe
portable power. Building the prototype from this parts palette means the design is not a
laboratory curiosity but a plausible product architecture that could be transitioned to a
custom board with modest additional effort.

Finally, the project aligns with the national and international policy emphasis on road
safety. Reducing deaths and injuries from road traffic accidents is an explicit target of
the United Nations Sustainable Development Goals, and India's own road-safety programmes
identify two-wheeler riders as the most vulnerable category of road user. A retrofit device
that shortens the time to help, and that costs less than a routine service, addresses that
priority at a price point at which it can realistically be adopted at scale.

% ============================================================
\chapter{LITERATURE REVIEW}

\section{SUMMARY OF EXISTING WORK}

Research on two-wheeler safety and ride telematics has developed along three largely
separate lines: satellite-based tracking of vehicle position, inertial detection of
accidents, and rider-assistance systems built on additional sensing. Each line has produced
capable results, but the literature contains relatively few attempts to combine all three
in a single low-cost unit, and almost none that do so without assuming continuous network
connectivity.

The tracking line of work is well established. A representative recent contribution is the
development of real-time GPS tracking with estimated time of arrival on an Android
application using an ESP32 and a u-blox NEO6MV2 receiver, presented at the IEEE
International Conference on Computing, Communication and Networking Technologies in 2024.
That work demonstrates convincingly that an inexpensive microcontroller paired with a
consumer-grade satellite receiver can acquire position reliably, compute derived quantities
such as arrival time, and deliver them to a mobile application with acceptable latency. It
confirms the viability of the microcontroller-plus-GNSS pairing adopted in the present
project. Its scope, however, stops at tracking: there is no inertial sensing, no accident
detection and no maintenance logic, and the design assumes a network path to the phone
rather than a direct local link.

A second line of work addresses rider safety through the Internet of Things more broadly.
Chhabra et al., writing in \textit{Scientific Reports} in 2024, propose an IoT-based smart
framework for the safe driving experience of two-wheelers, integrating sensing and alerting
into a connected architecture intended to improve rider outcomes. The contribution is
valuable in framing rider safety as a systems problem rather than a single-sensor problem,
and in demonstrating the benefit of coupling sensing to timely alerting. The framework
nevertheless depends on continuous network connectivity for its alerting path, which means
that its effectiveness degrades in exactly the low-coverage rural settings where a
two-wheeler crash is most likely to go unwitnessed.

A third line pursues higher-capability rider protection through additional sensing. Alai,
Rajamani and Sanjaya, in \textit{Mechanical Systems and Signal Processing} in 2024, describe
a smart e-scooter with embedded estimation of rear-vehicle trajectories for rider
protection. The work is technically sophisticated, showing that meaningful estimation of the
surrounding traffic can be performed on an embedded two-wheeler platform. Its requirements,
however, place it in a different cost class: ranging sensors and substantially greater
computational capability are needed, which makes the approach difficult to justify for a
mass-market retrofit unit intended to cost a small fraction of the vehicle's service bill.

Alongside these, a large body of academic prototypes implements accident alerting using
GSM or GPRS. The common pattern is an accelerometer that triggers above a fixed threshold,
whereupon a microcontroller composes an SMS containing GPS coordinates and sends it to a
stored contact number. These systems are simple, inexpensive in components and widely
reproduced, and they establish the basic feasibility of automatic crash notification.
Their weaknesses are equally well documented. They incur a recurring SIM cost; they fail
where cellular coverage is poor, which is disproportionately where help is most needed; and
because they trigger on a single acceleration threshold they are prone to false alarms from
potholes, speed breakers, kerb strikes and even the vehicle being knocked over while
parked. Repeated false alerts erode the user's confidence and, in practice, lead to the
device being disabled.

Two further strands of the literature inform the design choices made here. Work on
Bluetooth Low Energy for sensor-to-phone communication has established the GATT model as an
efficient, low-power means of exposing structured sensor data to a handset, with power
consumption low enough for battery-operated devices and range adequate for a device carried
on the same vehicle as its user. Separately, work on inertial fall and impact detection ---
much of it originating in wearable and elderly-care research --- has shown that combining
multiple conditions, typically an acceleration peak followed by an orientation change and a
period of inactivity, markedly reduces false positives compared with any single threshold.
This multi-condition principle, developed largely outside the vehicular context, is directly
applicable to two-wheeler crash detection and is adopted in the present work.

The table below summarises the principal works reviewed, their contributions and the gaps
they leave.

\vspace{0.6em}
\begin{center}
\renewcommand{\arraystretch}{1.2}
\begin{singlespace}
\begin{tabular}{|L{1.55in}|L{1.75in}|L{2.05in}|}
\hline
\textbf{Work} & \textbf{Contribution} & \textbf{Gap / Takeaway}\\
\hline
Real-Time GPS Tracking and ETA on an Android Application using ESP32 and u-blox NEO6MV2
(IEEE ICCCNT, 2024)
& Demonstrates ESP32 with a consumer GNSS receiver delivering position and estimated arrival
time to an Android application.
& Tracking only; no crash detection and no maintenance logic. Confirms the
microcontroller-plus-GNSS pairing adopted here.\\
\hline
Chhabra et al., \textit{IoT Based Smart Framework for the Safe Driving Experience of Two
Wheelers}, Scientific Reports (2024)
& IoT framework addressing rider safety through coupled sensing and alerting.
& Depends on continuous network connectivity; the present design keeps the critical path
fully offline over BLE.\\
\hline
Alai et al., \textit{A Smart E-Scooter with Embedded Estimation of Rear Vehicle
Trajectories for Rider Protection}, MSSP (2024)
& Embedded estimation of surrounding traffic on a two-wheeler platform.
& Requires ranging sensors and higher compute; too costly for a mass-market retrofit
unit.\\
\hline
Conventional GSM/GPRS accident-alert systems (widely reported academic prototypes)
& Accelerometer-triggered SMS alerts carrying GPS coordinates to a fixed contact.
& Recurring SIM cost, poor rural coverage and high false-alarm rates; motivates the
multi-condition crash logic used here.\\
\hline
BLE GATT sensor-to-phone literature
& Establishes GATT as a low-power, structured means of exposing sensor data to a handset.
& Validates BLE as the sole communication link, with the phone as client and the unit as
server.\\
\hline
Inertial fall-detection literature (wearable and elderly-care domain)
& Shows that combining impact, orientation change and post-event inactivity greatly reduces
false positives.
& Principle transfers directly to two-wheeler crash detection and is adopted in this
work.\\
\hline
\end{tabular}
\end{singlespace}
\captionof{table}{Summary of Existing Work and Identified Gaps}
\end{center}

\section{GAPS IDENTIFIED}

A review of the work above reveals consistent limitations that, taken together, explain why
no affordable two-wheeler ride-and-safety unit is in general use. The following gaps have
been identified:

\begin{itemize}[leftmargin=*,itemsep=0.6em,topsep=0.7em]
\item \textbf{Mandatory dependence on cellular connectivity.} Almost every safety-oriented
system reviewed assumes a GSM, GPRS or Wi-Fi path for its alerting function. This makes the
device useless in low-coverage areas, imposes a recurring data cost on the user, and
introduces a failure mode that is strongly correlated with the very locations where a crash
is most likely to go unnoticed.

\item \textbf{Single-threshold crash detection and consequent false alarms.} The prevailing
approach declares an accident whenever acceleration crosses a fixed limit. Ordinary riding
in Indian conditions routinely produces such peaks from potholes, speed breakers, kerb
strikes and abrupt braking. The resulting false alerts are the single most frequently
reported practical failure of these prototypes, and they lead directly to the device being
switched off.

\item \textbf{Separation of tracking, safety and maintenance functions.} Tracking systems
log position but do not detect crashes; accident-alert systems detect impacts but keep no
ride history; neither addresses servicing. The rider is therefore offered three partial
products rather than one coherent one, and no single reviewed system integrates all three.

\item \textbf{Absence of distance-based maintenance intelligence.} Although cumulative
distance is trivially available once position is being logged, none of the reviewed systems
uses it to drive servicing reminders. Maintenance consequently remains scheduled by
guesswork, with measurable cost to vehicle life and emissions.

\item \textbf{Cost and complexity of higher-capability approaches.} Systems that add ranging
sensors or significant computation achieve impressive rider protection but at a price and
power budget incompatible with a retrofit unit for the existing vehicle population.

\item \textbf{Cloud lock-in and loss of data ownership.} Commercial telematics routes the
rider's movement history through a proprietary back end. This raises privacy concerns over
continuous location data and leaves the device dependent on the continued operation of a
commercial service.

\item \textbf{No graceful behaviour when the link is unavailable.} Most reviewed designs
assume the communication channel is present. Few describe what happens to data recorded
while the channel is down, so records generated during a coverage gap are simply lost.

\item \textbf{Inadequate attention to power integrity during a crash.} Systems drawing power
directly from the vehicle supply are vulnerable to exactly the electrical disturbance a
crash produces. Few reviewed designs discuss how the unit stays alive long enough to report
the event it has just detected.

\item \textbf{Limited rider control over the alert.} Where an alert is raised, it is often
sent immediately and irrevocably. Without a cancellation window, a false positive
propagates to the emergency contact, compounding the trust problem described above.

\item \textbf{Sparse reporting of validation methodology.} Detection latency, false-alarm
rate and distance accuracy are seldom characterised systematically, which makes it difficult
to compare approaches or to establish that thresholds generalise beyond the authors' own
trials.
\end{itemize}

Taken together, these gaps point to the need for a low-cost, self-contained on-board unit
that integrates ride logging, multi-condition crash detection and distance-based maintenance
in one device, that keeps its critical path entirely offline, that buffers data when the
link is down, that remains powered through an impact, and that leaves the rider in control
of the alert.

\section{CONTRIBUTION OF THE PROPOSED WORK}

The proposed project, \textit{SmartRide: Bluetooth-Based Two-Wheeler Ride Logger and Crash
Detection System Using ESP32-S3}, addresses the identified gaps through the following
contributions:

\begin{itemize}[leftmargin=*,itemsep=0.6em,topsep=0.7em]
\item \textbf{A fully offline, network-independent architecture.} SmartRide eliminates the
cellular link entirely. Sensing, filtering, ride computation, crash reasoning and storage
all take place on the on-board unit and the paired handset, connected by a short Bluetooth
Low Energy link. The system therefore functions identically on a rural road and in a city,
requires no SIM card or data plan, and incurs no recurring cost.

\item \textbf{Multi-condition crash detection that suppresses false alarms.} Instead of a
single acceleration threshold, the firmware requires three conditions in sequence: a
resultant-acceleration peak above a calibrated limit, a sustained tilt beyond an
orientation bound derived from gyroscope integration, and a period of near-stillness
following the impact. Because a pothole produces an acceleration spike without a sustained
tilt, and a low-speed manoeuvre produces tilt without an impact, the conjunction of all
three rejects the dominant false-alarm sources reported in the literature.

\item \textbf{An optional redundant impact channel.} The ADXL345 may be populated alongside
the MPU6500 on the same bus at a distinct address. Its built-in tap and activity engines
raise a hardware interrupt independently of the main sampling loop, providing corroboration
of a high-energy event from a physically separate device. This addresses sensor-level
single-point failure at negligible cost, and the firmware degrades gracefully to
MPU6500-only operation when the part is not fitted.

\item \textbf{Integration of ride logging, safety and maintenance in one unit.} A single
device provides the ride record, the crash alert and the servicing reminder. Cumulative
distance, already computed for the ride log, is reused to drive maintenance reminders at
user-configurable intervals, so the maintenance feature is obtained at essentially no
additional hardware cost.

\item \textbf{Edge computation on the ESP32-S3 with the phone as a pure client.} All ride
statistics are computed on the microcontroller and exposed as BLE GATT characteristics.
The application reads and renders them but performs no essential computation, which means
the unit keeps logging correctly whether or not the phone is connected, and the same unit
would work unchanged against a future iOS client.

\item \textbf{Local buffering with deferred delivery.} Records and crash events generated
while the handset is absent or out of range are written to non-volatile storage on the unit
and transferred on the next successful connection. No data is lost to a link outage, which
directly addresses one of the gaps identified above.

\item \textbf{Independent, crash-survivable power.} The unit runs from its own lithium-ion
cell managed by a dedicated battery management system providing over-charge,
over-discharge, over-current and short-circuit protection. Because the logger does not
depend on the vehicle supply at the moment of impact, it remains alive to report the event
it has just detected.

\item \textbf{Multi-constellation positioning suited to Indian conditions.} The
ATGM336H-5N receives from several satellite constellations concurrently, which improves fix
availability and the time to first fix relative to a single-constellation receiver under
obstructed sky. Its one-pulse-per-second output is used to discipline record timestamps so
that logged data remains temporally consistent.

\item \textbf{A rider-cancellable alert.} A crash event raises a countdown in the
application which the rider can dismiss. This keeps the human in the loop, ensures that a
residual false positive does not propagate, and preserves the rider's trust in the device
--- the practical precondition for it remaining switched on.

\item \textbf{Privacy by architecture.} Because no movement data leaves the vehicle and the
rider's own handset, location privacy is a structural property of the design rather than a
policy promise, and the device cannot be rendered inoperable by the discontinuation of a
commercial service.

\item \textbf{A defined validation methodology.} Distance and speed accuracy, crash
detection latency, false-alarm behaviour under deliberately adverse road inputs, and BLE
link reliability over realistic rider-to-phone separation are each to be characterised
explicitly, addressing the weak validation reporting observed in the reviewed literature.
\end{itemize}
% ============================================================
\chapter{METHODOLOGY}

\section{APPROACH}

The proposed system, \textit{SmartRide}, follows a hybrid experimental and analytical
methodology in which hardware bring-up, firmware development and mobile application
development proceed in parallel and are brought together through incremental integration.
The guiding principle is that every function on the critical path must execute either on the
on-board unit or on the paired handset, so that the system never depends on a network to do
its job.

Work begins with sensor bring-up on a breadboard. The ATGM336H-5N GNSS receiver is connected
to a hardware UART of the ESP32-S3 and its NMEA output is parsed and validated against known
coordinates; the MPU6500 inertial measurement unit is connected over its serial interface and
its accelerometer and gyroscope scales are verified against static gravity and controlled
rotation. Where the optional ADXL345 is populated, it is placed on the same
I\textsuperscript{2}C bus at its alternate address and its interrupt output is routed to a
free GPIO. This stage establishes, for each device, a known-good register configuration, a
sampling rate and a noise characterisation that the later stages depend upon.

The firmware is then developed as a set of cooperating real-time tasks rather than a single
polling loop. A high-rate task services the inertial measurement unit through its on-chip
FIFO so that no samples are lost to scheduling jitter; a low-rate task parses GNSS sentences
as they arrive; a computation task consumes both streams from queues and derives ride
statistics; and a communication task maintains the Bluetooth Low Energy connection and
publishes updated characteristic values. Separating these concerns means that a slow or
absent phone connection cannot delay sensor acquisition, and that the crash rule is evaluated
on a deterministic schedule.

Ride statistics are computed on the microcontroller itself. Instantaneous speed is taken
from the GNSS velocity solution rather than differentiated from position, which is
substantially less noisy. Incremental distance between consecutive accepted fixes is
computed with the Haversine formula and accumulated; fixes are accepted only when the
receiver reports a valid solution and the implied speed is physically plausible, which
prevents the cumulative odometer from being corrupted by outliers. Ride duration is taken
from timestamps disciplined by the receiver's one-pulse-per-second output so that records
remain temporally consistent across a long ride.

Crash detection is treated as a sequential decision rather than a single comparison. The
firmware first looks for an impact in the resultant acceleration, then requires the vehicle's
orientation to have changed and remained changed, and finally requires a period of
near-stillness. Only when all three conditions have occurred in that order within a bounded
time window is a crash declared. The design intent is explicitly to accept a slightly
longer detection latency in exchange for a false-alarm rate low enough that the rider leaves
the device switched on.

The mobile application is developed against the unit's BLE GATT service using Flutter. It
subscribes to the ride and crash characteristics, renders a live dashboard, plots the route
on a map, and writes each completed ride into a local SQLite database from which history and
the cumulative odometer are served. A crash notification raises a countdown that the rider
can cancel, keeping the human in the loop.

Development is iterative. Each stage is validated before the next is layered on: sensor
readings against physical ground truth, distance against known routes, the crash rule
against controlled drop and tilt trials and against deliberately adverse road inputs such as
speed breakers and kerb strikes, and the BLE link against realistic rider-to-phone
separation. This closed-loop cycle --- sense, compute, decide, deliver, measure, retune ---
is what converts a set of working components into a system whose behaviour can be trusted.

\begin{figure}[htbp]
\centering
\begin{tikzpicture}[
  font=\sffamily\scriptsize,
  box/.style={draw,rounded corners=2pt,align=center,inner sep=4pt,minimum height=9mm,
              text width=25mm},
  core/.style={draw,rounded corners=2pt,align=center,inner sep=4pt,minimum height=16mm,
               text width=31mm,fill=blue!10,thick},
  app/.style={draw,rounded corners=2pt,align=center,inner sep=4pt,minimum height=11mm,
              text width=29mm,fill=gray!8},
  ar/.style={-{Latex[length=2mm]},thick}
]
% left column : sensors and power
\node[box] (bat)  at (0,2.15)  {Li-ion Cell + BMS\\(\SI{3.7}{V} regulated)};
\node[box] (gnss) at (0,0.75)  {ATGM336H-5N GNSS\\(UART + 1\,PPS)};
\node[box] (imu)  at (0,-0.65) {MPU6500 IMU\\(SPI / I\textsuperscript{2}C)};
\node[box,dashed] (adxl) at (0,-2.05) {ADXL345\\(optional, I\textsuperscript{2}C + INT)};

% core
\node[core] (esp) at (4.4,0.05) {\textbf{ESP32-S3}\\Sampling, filtering,\\ride statistics,\\crash rule, BLE server};

% right column
\node[app] (sql) at (9.3,1.85) {SQLite\\Ride history};
\node[app] (flut) at (9.3,0.05) {Flutter Android App\\BLE client};
\node[app] (maps) at (9.3,-1.75) {Google Maps API\\Route display};

\draw[ar] (bat)  -- (esp);
\draw[ar] (gnss) -- (esp);
\draw[ar] (imu)  -- (esp);
\draw[ar,dashed] (adxl) -- (esp);
\draw[ar] (esp)  -- node[above,font=\sffamily\tiny,align=center]{BLE}
                   node[below,font=\sffamily\tiny,align=center]{GATT} (flut);
\draw[ar] (flut) -- (sql);
\draw[ar] (sql)  -- (flut);
\draw[ar] (flut) -- (maps);
\draw[ar] (maps) -- (flut);
\end{tikzpicture}
\caption{System Architecture of SmartRide}
\end{figure}

\section{THEORETICAL FRAMEWORK}

The theoretical foundation of SmartRide rests on three bodies of principle: the geodesy used
to convert satellite fixes into distance travelled, the mechanics used to recognise an
impact and a fall from inertial data, and the communication model used to expose the results
to a handset.

\subsection{Positioning, Speed and Distance}

The GNSS receiver reports position as a latitude and longitude pair in the WGS-84 datum.
The great-circle distance between two consecutive fixes
\((\varphi_1,\lambda_1)\) and \((\varphi_2,\lambda_2)\) is obtained from the Haversine
formula:

\begin{equation}
a = \sin^{2}\!\left(\frac{\Delta\varphi}{2}\right)
  + \cos\varphi_1\,\cos\varphi_2\,\sin^{2}\!\left(\frac{\Delta\lambda}{2}\right)
\end{equation}

\begin{equation}
d = 2R\,\operatorname{atan2}\!\left(\sqrt{a},\,\sqrt{1-a}\right)
\end{equation}

where \(\Delta\varphi = \varphi_2-\varphi_1\), \(\Delta\lambda = \lambda_2-\lambda_1\), and
\(R\) is the mean Earth radius, taken as \SI{6371}{km}. The formula is preferred over the
simpler equirectangular approximation because it remains accurate for the short baselines
between successive fixes without special handling near the poles, and it is inexpensive
enough to evaluate on the microcontroller at the receiver's update rate.

The cumulative ride distance is the sum of accepted increments:

\begin{equation}
D = \sum_{i=1}^{N} d_i , \qquad
d_i \text{ accepted iff } \text{fix valid} \ \wedge\ \frac{d_i}{\Delta t_i} \le v_{\max}
\end{equation}

The plausibility gate \(d_i/\Delta t_i \le v_{\max}\) is essential in practice: a single
spurious fix would otherwise inject a large spurious increment into the odometer that could
never be removed. Instantaneous speed is taken directly from the receiver's velocity
solution, which is derived from Doppler measurements and is markedly less noisy than a
position difference:

\begin{equation}
v = v_{\text{GNSS}} , \qquad
\bar v = \frac{D}{T} , \qquad T = \sum_{i} \Delta t_i
\end{equation}

\subsection{Impact and Orientation Mechanics}

The MPU6500 reports acceleration along three orthogonal body axes. The resultant, or
magnitude, of the acceleration vector is

\begin{equation}
a_{r} = \sqrt{a_x^{2} + a_y^{2} + a_z^{2}}
\end{equation}

At rest this quantity equals \(1\,g\) regardless of the orientation of the unit, which makes
it a convenient orientation-independent impact indicator. An impact is flagged when
\(a_r\) exceeds a calibrated threshold \(a_{\text{th}}\), nominally in the region of
\(4\,g\) and tuned empirically during field trials.

Orientation is obtained by integrating the gyroscope's angular rates and correcting the
result against the accelerometer's estimate of the gravity direction, which bounds the drift
inherent in integration alone. The tilt of the unit from the vertical follows from the
accelerometer directly:

\begin{equation}
\theta = \arccos\!\left(\frac{a_z}{a_{r}}\right)
\end{equation}

A fall is characterised not by a momentary tilt --- a motorcycle leans substantially in
normal cornering --- but by a tilt that exceeds a bound and is then \emph{sustained}:

\begin{equation}
\theta(t) > \theta_{\text{th}} \quad \text{for all } t \in [t_0,\, t_0 + \tau_{\text{tilt}}]
\end{equation}

with \(\theta_{\text{th}}\) nominally \(60^{\circ}\). The final condition is stillness:
over a short window following the impact, the variation of the resultant acceleration must
remain small,

\begin{equation}
\max_{t \in [t_1,\, t_1+\tau_{\text{still}}]} \bigl| a_{r}(t) - 1g \bigr|
  < \epsilon_{\text{still}}
\end{equation}

The crash decision is the conjunction of the three, evaluated in order within a bounded
window:

\begin{equation}
\mathrm{Crash} = \bigl(a_r > a_{\text{th}}\bigr)
  \ \wedge\ \bigl(\theta \text{ sustained} > \theta_{\text{th}}\bigr)
  \ \wedge\ \bigl(\text{stillness}\bigr)
\end{equation}

The value of the conjunction lies in what each term excludes. A pothole or speed breaker
produces a large \(a_r\) but no sustained tilt, because the vehicle remains upright and
continues moving. A slow-speed lean or a tight turn produces tilt without an impact. A
vehicle stopped at a signal satisfies stillness alone. Only a genuine fall produces all
three in sequence, which is why the multi-condition rule is markedly more selective than any
single threshold.

\subsection{Communication Model}

Bluetooth Low Energy structures data exchange through the Generic Attribute Profile (GATT).
The on-board unit acts as a GATT \emph{server}, exposing a custom primary service whose
characteristics carry ride statistics, route points and crash events; the handset acts
purely as a \emph{client}, subscribing to notifications and reading values on demand. This
assignment of roles is deliberate. Because the unit is the server and the source of truth,
it continues to sample sensors, compute statistics and evaluate the crash rule whether or
not a client is connected, and records generated while the phone is absent are written to
non-volatile storage and delivered on the next connection. The phone performs no essential
computation, so the same firmware would serve an iOS or desktop client without
modification.

\section{TOOLS AND TECHNOLOGIES}

\subsection{Hardware Components}

\begin{itemize}[leftmargin=*,itemsep=0.6em,topsep=0.6em]
\item \textbf{ESP32-S3 Microcontroller.} A dual-core Xtensa LX7 system on chip running at up
to \SI{240}{MHz}, with integrated \SI{2.4}{GHz} Wi-Fi and Bluetooth~5 Low Energy, \SI{512}{kB}
of on-chip SRAM, optional external PSRAM and a rich complement of serial peripherals. It
acts as the sole processor of the on-board unit: it reads the GNSS receiver over UART and the
inertial measurement unit over SPI or I\textsuperscript{2}C, computes ride statistics,
evaluates the crash rule, maintains the non-volatile record buffer and hosts the BLE GATT
service. The part was chosen over the original ESP32 for its greater processing headroom,
its native USB interface, which simplifies flashing and log capture, and its Bluetooth~5 Low
Energy radio.

\item \textbf{ATGM336H-5N GNSS Receiver.} A compact multi-constellation satellite receiver
that tracks several systems concurrently and outputs standard NMEA~0183 sentences over UART,
together with a one-pulse-per-second timing signal. Concurrent multi-constellation
reception improves fix availability and shortens time-to-first-fix under obstructed sky
compared with a single-constellation receiver, which matters on tree-lined and built-up
Kerala roads. The module supplies latitude, longitude, ground speed, course and satellite
count, and its 1\,PPS output is used to discipline record timestamps.

\item \textbf{MPU6500 Inertial Measurement Unit.} A six-axis device combining a three-axis
accelerometer and a three-axis gyroscope with on-chip sixteen-bit conversion, selectable
full-scale ranges up to $\pm 16\,g$ and $\pm 2000^{\circ}/\mathrm{s}$, a configurable digital
low-pass filter, an internal FIFO and a wake-on-motion interrupt. It supplies the
acceleration and angular-rate streams from which impact and sustained tilt are derived.
Reading through the FIFO rather than polling individual registers ensures that no samples
are lost to scheduling jitter in the firmware.

\item \textbf{ADXL345 Accelerometer (Optional).} A three-axis digital accelerometer with
selectable ranges up to $\pm 16\,g$, thirteen-bit resolution and hardware detection engines
for single tap, double tap, activity, inactivity and free fall, any of which can raise an
interrupt on a dedicated pin. When populated, it sits on the same I\textsuperscript{2}C bus
at its alternate address and provides an independent, physically separate corroboration of a
high-energy event without consuming processor time. The firmware treats its presence as
optional and degrades gracefully to MPU6500-only operation when it is absent.

\item \textbf{Lithium-Ion Cell with Battery Management System.} A single Li-ion cell,
already available with the team, supplies the unit through a dedicated battery management
system providing over-charge, over-discharge, over-current and short-circuit protection,
with a regulator presenting a clean \SI{3.3}{V} rail to the logic. Running the unit from its
own managed cell rather than directly from the vehicle supply is a deliberate design
decision: the logger must survive the electrical disturbance of the very event it is
designed to report.

\item \textbf{Enclosure, Mounting and Wiring.} A closed enclosure with vibration-isolating
mounting protects the assembly and, importantly, fixes the orientation of the inertial
measurement unit relative to the vehicle, which the tilt calculation depends upon. The GNSS
antenna is positioned for clear sky visibility.
\end{itemize}

\subsection{Software Tools}

\begin{itemize}[leftmargin=*,itemsep=0.55em,topsep=0.6em]
\item \textbf{ESP-IDF and the Arduino Core for ESP32-S3.} The primary firmware development
environment, providing the FreeRTOS kernel, peripheral drivers, the BLE host stack and the
non-volatile storage subsystem used for the record buffer.

\item \textbf{Flutter (Dart).} Used to build the Android application, giving a single
codebase for the dashboard, map view, ride history and reminder logic, with mature plugins
for Bluetooth Low Energy and mapping.

\item \textbf{SQLite.} The embedded relational database within the application, holding
completed rides, route point sequences, crash events and the cumulative odometer from which
maintenance reminders are derived.

\item \textbf{Google Maps API.} Used within the application to render the travelled route
from the logged coordinate sequence.

\item \textbf{Python with NumPy and Matplotlib.} Used offline to analyse captured sensor
logs, to visualise impact and tilt signatures from controlled trials, and to tune the crash
thresholds before they are committed to firmware.

\item \textbf{Git and GitHub.} Version control for firmware, application and documentation,
with the repository also serving as the record of design decisions across the phases.

\item \textbf{Programming Languages.} C and C++ for the firmware, Dart for the mobile
application, and Python for analysis.
\end{itemize}

\section{DESIGN PROCESS}

The design is presented through a use case diagram that fixes the boundary of the system and
the interactions across it, a hierarchy of data flow diagrams that decompose the system
first into its major components and then into the processing performed within the
microcontroller, and a sequence diagram that establishes the temporal ordering of a ride and
of a crash event. Together these views describe a system involving a Rider, an on-board unit
comprising the ESP32-S3 with the ATGM336H-5N GNSS receiver and the MPU6500 inertial
measurement unit, and a Flutter mobile application.

\subsection{Use Case Diagram for the Ride Logging and Crash Detection System}

The interactions fall into three groups: ride operation, safety, and history and
maintenance.

\begin{figure}[htbp]
\centering
\begin{tikzpicture}[
  font=\sffamily\scriptsize,
  uc/.style={draw,ellipse,align=center,inner sep=2pt,minimum width=34mm,minimum height=8mm},
  ln/.style={-,thin}
]
\draw (2.55,-5.15) rectangle (8.05,0.55);
\node[font=\sffamily\scriptsize] at (5.3,0.28) {\textbf{SmartRide System}};

\node[uc] (u1) at (5.3,-0.35) {Pair On-Board Unit};
\node[uc] (u2) at (5.3,-1.05) {Start / Stop Ride};
\node[uc] (u3) at (5.3,-1.75) {View Live Dashboard};
\node[uc] (u4) at (5.3,-2.45) {View Route on Map};
\node[uc] (u5) at (5.3,-3.15) {View Ride History};
\node[uc] (u6) at (5.3,-3.85) {Receive Maintenance\\Reminder};
\node[uc] (u7) at (5.3,-4.65) {Cancel False Crash Alert};

\node[uc] (s1) at (5.3,-6.05) {Log Ride Data};
\node[uc] (s2) at (5.3,-6.85) {Detect Crash and\\Raise Alert};
\draw (2.55,-7.35) rectangle (8.05,-5.6);

% Rider actor
\begin{scope}[shift={(0.55,-2.5)}]
  \draw (0,0.32) circle (0.13);
  \draw (0,0.19) -- (0,-0.22);
  \draw (-0.22,0.02) -- (0.22,0.02);
  \draw (0,-0.22) -- (-0.17,-0.55);
  \draw (0,-0.22) -- (0.17,-0.55);
  \node[font=\sffamily\scriptsize,below] at (0,-0.62) {Rider};
\end{scope}

% On-board unit actor
\begin{scope}[shift={(9.95,-6.45)}]
  \draw (0,0.32) circle (0.13);
  \draw (0,0.19) -- (0,-0.22);
  \draw (-0.22,0.02) -- (0.22,0.02);
  \draw (0,-0.22) -- (-0.17,-0.55);
  \draw (0,-0.22) -- (0.17,-0.55);
  \node[font=\sffamily\scriptsize,below,align=center] at (0,-0.62) {On-Board\\Unit};
\end{scope}

\foreach \t in {u1,u2,u3,u4,u5,u6,u7} { \draw[ln] (0.78,-2.5) -- (\t); }
\draw[ln] (0.78,-2.5) -- (s2);
\draw[ln] (9.7,-6.45) -- (s1);
\draw[ln] (9.7,-6.45) -- (s2);
\end{tikzpicture}
\caption{Use Case Diagram of the SmartRide System}
\end{figure}

\vspace{0.4em}
\noindent\textbf{(i) Ride Operation.} The rider pairs the on-board unit with the handset
once, after which the application reconnects automatically whenever the unit is in range.

\begin{itemize}[leftmargin=*,itemsep=0.3em,topsep=0.45em]
\item \textbf{Pair On-Board Unit:} a one-time Bluetooth Low Energy pairing that binds the
unit to the rider's handset.
\item \textbf{Start / Stop Ride:} delimits a ride so that statistics and route points are
grouped into a single record. The unit also infers ride boundaries from sustained motion
and stillness, so an explicit action is not required.
\item \textbf{View Live Dashboard:} presents instantaneous speed, elapsed distance, ride
duration and satellite status as the ride proceeds.
\item \textbf{View Route on Map:} plots the logged coordinate sequence over a map view.
\end{itemize}

\noindent\textbf{(ii) Safety.} The safety interactions are the ones the rider hopes never to
use.

\begin{itemize}[leftmargin=*,itemsep=0.3em,topsep=0.45em]
\item \textbf{Detect Crash and Raise Alert:} performed by the on-board unit, which evaluates
the multi-condition rule continuously and notifies the application when all conditions are
met.
\item \textbf{Cancel False Crash Alert:} the rider dismisses the countdown raised by the
application, keeping a residual false positive from propagating further.
\end{itemize}

\noindent\textbf{(iii) History and Maintenance.}

\begin{itemize}[leftmargin=*,itemsep=0.3em,topsep=0.45em]
\item \textbf{Log Ride Data:} performed by the on-board unit, which writes records to its
non-volatile buffer and transfers them on connection.
\item \textbf{View Ride History:} browses past rides held in the handset's SQLite database.
\item \textbf{Receive Maintenance Reminder:} raised by the application when the cumulative
odometer crosses a user-configured servicing interval.
\end{itemize}

\subsection{System Overview: SmartRide Data Flow (Level 0 DFD)}

The context diagram identifies what lies outside the system and what crosses its boundary.
Three external entities are relevant. The \textbf{GNSS constellation} supplies the satellite
signals from which position, speed and course are derived. The \textbf{vehicle and road}
constitute the physical environment whose motion the inertial sensors observe. The
\textbf{rider} supplies configuration and control and consumes the system's output.

\begin{figure}[htbp]
\centering
\begin{tikzpicture}[
  font=\sffamily\scriptsize,
  ent/.style={draw,align=center,inner sep=4pt,minimum height=9mm,text width=24mm},
  proc/.style={draw,circle,align=center,inner sep=2pt,minimum size=30mm,fill=blue!8,thick},
  ar/.style={-{Latex[length=2mm]},thick},
  lbl/.style={font=\sffamily\tiny,align=center,fill=white,inner sep=1.5pt}
]
\node[ent] (sat)   at (0,2.3)   {GNSS\\Constellation};
\node[ent] (road)  at (0,-2.3)  {Vehicle and Road\\(motion, impact)};
\node[proc] (sys)  at (6.0,0)   {\textbf{0}\\SmartRide\\Ride Logging and\\Crash Detection\\System};
\node[ent] (rider) at (11.4,0)  {Rider};

\draw[ar] (sat)  -- node[lbl,sloped,pos=0.45,above]
                    {position, speed,\\course, 1\,PPS} (sys);
\draw[ar] (road) -- node[lbl,sloped,pos=0.45,below]
                    {acceleration,\\angular rate} (sys);
\draw[ar] (sys) to[bend left=16] node[lbl,pos=0.5,above]
                    {ride statistics, route,\\crash alert, reminder} (rider);
\draw[ar] (rider) to[bend left=16] node[lbl,pos=0.5,below]
                    {start/stop, service interval,\\alert cancellation} (sys);
\end{tikzpicture}
\caption{Level 0 DFD: SmartRide Context Diagram}
\end{figure}

The system comprises four core elements: the ATGM336H-5N GNSS receiver, the MPU6500 inertial
measurement unit, the ESP32-S3 microcontroller and the Flutter mobile application. Data
flows from the sensing components through the microcontroller to the application, with a
control path returning from the rider.

\vspace{0.5em}
\noindent\textbf{(i) ATGM336H-5N --- GNSS Receiver.} The module acquires and tracks signals
from multiple satellite constellations concurrently and computes a navigation solution.

\noindent\emph{Process:}
\begin{itemize}[leftmargin=*,itemsep=0.28em,topsep=0.4em]
\item It acquires satellite signals and resolves a position, velocity and time solution.
\item It emits the solution as NMEA~0183 sentences over UART at a configured update rate.
\item It asserts a one-pulse-per-second signal aligned to satellite time.
\end{itemize}
\noindent\emph{Output:} latitude, longitude, ground speed, course, fix validity and
satellite count, delivered to the ESP32-S3 over UART, with 1\,PPS on a separate GPIO for
timestamp discipline.

\vspace{0.5em}
\noindent\textbf{(ii) MPU6500 --- Inertial Measurement Unit.} The device measures linear
acceleration and angular rate along three orthogonal axes.

\noindent\emph{Process:}
\begin{itemize}[leftmargin=*,itemsep=0.28em,topsep=0.4em]
\item Micro-machined structures respond to acceleration and rotation, and the responses are
digitised on chip at sixteen-bit resolution.
\item A configurable digital low-pass filter removes out-of-band content before the samples
reach the FIFO.
\item Samples accumulate in the internal FIFO, from which the firmware reads them in blocks.
\end{itemize}
\noindent\emph{Output:} time-ordered acceleration and angular-rate samples delivered to the
ESP32-S3, from which the resultant acceleration and the tilt angle are computed. Where the
ADXL345 is populated, it additionally asserts a hardware interrupt on a high-energy event.

\vspace{0.5em}
\noindent\textbf{(iii) ESP32-S3 --- Microcontroller.} The microcontroller is the entire
intelligence of the on-board unit.

\noindent\emph{Process:}
\begin{itemize}[leftmargin=*,itemsep=0.28em,topsep=0.4em]
\item \textbf{Acquisition:} it parses NMEA sentences as they arrive and drains the inertial
FIFO on a deterministic schedule.
\item \textbf{Computation:} it validates fixes, computes Haversine increments, accumulates
distance and duration, and maintains the ride record.
\item \textbf{Decision:} it evaluates the sequential crash rule over the inertial stream.
\item \textbf{Persistence:} it writes records to non-volatile storage whenever no client is
connected.
\item \textbf{Publication:} it hosts the BLE GATT service and notifies subscribed clients of
updated values and crash events.
\end{itemize}
\noindent\emph{Output:} ride statistics, route points and crash events, published over
Bluetooth Low Energy.

\vspace{0.5em}
\noindent\textbf{(iv) Flutter Mobile Application.} The application is the presentation and
archival layer.

\noindent\emph{Process:} it connects as a BLE client, subscribes to the unit's
characteristics, renders the dashboard and map, writes completed rides to SQLite, maintains
the cumulative odometer and raises maintenance reminders and cancellable crash alerts.

\noindent\emph{Output:} the rider's view of the system, and the durable ride archive.

\subsection{System Data Management and Delivery (Level 1 DFD)}

The Level 1 decomposition expands the single process of the context diagram into the
processes that actually run, and names the two stores in which data persists.

\begin{figure}[htbp]
\centering
\begin{tikzpicture}[
  font=\sffamily\scriptsize,
  ent/.style={draw,align=center,inner sep=3pt,minimum height=8mm,text width=21mm},
  proc/.style={draw,rounded corners=8pt,align=center,inner sep=3pt,minimum height=11mm,
               text width=26mm,fill=blue!8},
  store/.style={draw,align=center,inner sep=3pt,minimum height=9mm,text width=26mm,
                fill=gray!12},
  ar/.style={-{Latex[length=2mm]},thick},
  lbl/.style={font=\sffamily\tiny,align=center,fill=white,inner sep=1.5pt}
]
\node[ent]   (sat)  at (0,3.3)    {GNSS\\Constellation};
\node[proc]  (p1)   at (3.9,3.3)  {\textbf{1.0}\\Acquire and Parse\\GNSS Data};
\node[proc]  (p3)   at (8.1,3.3)  {\textbf{3.0}\\Compute Ride\\Statistics};

\node[ent]   (road) at (0,1.4)    {Vehicle\\and Road};
\node[proc]  (p2)   at (3.9,1.4)  {\textbf{2.0}\\Acquire and Filter\\Motion Data};
\node[proc]  (p4)   at (8.1,1.4)  {\textbf{4.0}\\Evaluate Crash\\Rule};

\node[store] (d1)   at (3.6,-0.7) {\textbf{D1}\\Flash Record Buffer (NVS)};
\node[proc]  (p5)   at (8.1,-0.7) {\textbf{5.0}\\BLE GATT\\Server};

\node[store] (d2)   at (3.6,-3.0) {\textbf{D2}\\SQLite Ride History};
\node[proc]  (p6)   at (8.1,-3.0) {\textbf{6.0}\\Mobile App\\Presentation};

\node[ent]   (rider) at (8.1,-5.0) {Rider};

\draw[ar] (sat)  -- (p1);
\draw[ar] (road) -- (p2);
\draw[ar] (p1)   -- node[lbl,above]{fixes} (p3);
\draw[ar] (p2)   -- node[lbl,above]{$a_x,a_y,a_z,\omega$} (p4);
\draw[ar] (p1)   -- node[lbl,right]{1\,PPS} (p2);

% 3.0 -> 5.0 routed around the right
\draw[ar] (p3.east) -- ++(1.35,0) |- node[lbl,pos=0.28,right]{statistics} (p5.east);
\draw[ar] (p4)   -- node[lbl,right]{crash event} (p5);

% 5.0 <-> D1
\draw[ar] ([yshift=2.2mm]p5.west) -- node[lbl,above]{store when offline}
          ([yshift=2.2mm]d1.east);
\draw[ar] ([yshift=-2.2mm]d1.east) -- node[lbl,below]{drain on connect}
          ([yshift=-2.2mm]p5.west);

\draw[ar] (p5) -- node[lbl,right]{BLE notify} (p6);

% 6.0 <-> D2
\draw[ar] ([yshift=2.2mm]p6.west) -- node[lbl,above]{completed rides}
          ([yshift=2.2mm]d2.east);
\draw[ar] ([yshift=-2.2mm]d2.east) -- node[lbl,below]{history, odometer}
          ([yshift=-2.2mm]p6.west);

% 6.0 <-> Rider
\draw[ar] ([xshift=-7mm]p6.south) -- node[lbl,left,align=right]
          {dashboard, map,\\alert, reminder} ([xshift=-7mm]rider.north);
\draw[ar] ([xshift=7mm]rider.north) -- node[lbl,right,align=left]
          {start/stop,\\cancel} ([xshift=7mm]p6.south);

% control relayed back to the unit
\draw[ar,dashed] (p6.east) -- ++(2.3,0) |- node[lbl,pos=0.26,right]{relayed} (p5.south east);
\end{tikzpicture}
\caption{Level 1 DFD: SmartRide Data Management and Delivery}
\end{figure}

\vspace{0.5em}
\noindent\textbf{(i) Data Ingestion and Processing}

Process~1.0 reads bytes from the UART attached to the GNSS receiver, reassembles them into
complete NMEA sentences, verifies each checksum and extracts latitude, longitude, ground
speed, course, fix validity and satellite count. Sentences failing the checksum are
discarded silently; fixes marked invalid are not allowed to reach the distance computation.

Process~2.0 drains the inertial measurement unit's FIFO in blocks, converts the raw
sixteen-bit counts to physical units using the configured full-scale ranges, applies bias
corrections established during calibration, and smooths the stream with a short moving
average that removes sensor noise without masking the sharp transient of an impact. Where
the optional ADXL345 is fitted, its interrupt line is monitored here as an independent
high-energy indication.

Process~3.0 computes the ride statistics. It applies the plausibility gate to each candidate
fix, evaluates the Haversine increment, accumulates distance and duration, and appends the
coordinate to the route sequence. Process~4.0 evaluates the sequential crash rule --- impact,
sustained tilt, stillness --- over the filtered inertial stream, maintaining the small amount
of state required to recognise the conditions in order within a bounded window.

\vspace{0.5em}
\noindent\textbf{(ii) Persistence and Publication Paths}

Output from the computation processes follows two paths. \textbf{Store D1}, a non-volatile
flash region on the microcontroller, receives ride records whenever no client is connected,
so that a ride taken without the handset present is not lost. \textbf{Process 5.0}, the BLE
GATT server, publishes live statistics and crash events to a connected client, and on
connection first drains any records held in D1 before resuming live notification. This
ordering matters: the rider's history remains complete and in sequence regardless of when
the phone was in range.

\vspace{0.5em}
\noindent\textbf{(iii) Alert and User Delivery}

\textbf{Process 6.0}, the mobile application, is the final interface. It writes completed
rides and their route sequences into \textbf{Store D2}, the SQLite database on the handset,
from which ride history is browsed and the cumulative odometer is maintained. It compares
the odometer against the configured servicing interval and raises a maintenance reminder
when the interval is crossed. On receiving a crash event it raises an alert with a
countdown, which the rider may cancel; the cancellation is recorded and returned to the
unit, both to stop the alert and to inform later threshold review.

In summary, the two sensing processes feed two computation processes, which branch into a
persistence path through the flash buffer and a live publication path through the GATT
server, converging at the mobile application where data is archived and presented to the
rider.

\subsection{Data Processing Flow within the ESP32-S3 Microcontroller}

This diagram expands the processing internal to the microcontroller, showing how raw sensor
input becomes a published record. The firmware is organised as a set of FreeRTOS tasks
communicating through queues, so that acquisition, computation and communication proceed
without blocking one another.

\begin{figure}[htbp]
\centering
\begin{tikzpicture}[
  font=\sffamily\scriptsize,
  src/.style={draw,align=center,inner sep=3pt,minimum height=8mm,text width=23mm,
             fill=gray!8},
  tk/.style={draw,rounded corners=3pt,align=center,inner sep=3pt,minimum height=9mm,
             text width=27mm,fill=blue!8},
  q/.style={draw,align=center,inner sep=2pt,minimum height=6mm,text width=17mm,
            fill=yellow!15},
  snk/.style={draw,rounded corners=3pt,align=center,inner sep=3pt,minimum height=9mm,
              text width=24mm,fill=green!8},
  ar/.style={-{Latex[length=2mm]},thick}
]
% inputs
\node[src] (gin)  at (0,2.2)   {ATGM336H-5N\\UART + 1\,PPS};
\node[src] (iin)  at (0,0.0)   {MPU6500\\FIFO (SPI)};
\node[src,dashed] (ain) at (0,-2.0) {ADXL345 INT\\(optional)};

% acquisition tasks
\node[tk] (t1) at (3.5,2.2)  {\textbf{gnss\_task}\\sentence reassembly,\\checksum, parse};
\node[tk] (t2) at (3.5,0.0)  {\textbf{imu\_task}\\FIFO drain, scaling,\\bias, moving average};

% queues
\node[q] (q1) at (6.6,2.2)  {fix queue};
\node[q] (q2) at (6.6,0.0)  {motion queue};

% compute
\node[tk] (t3) at (9.6,2.4)  {\textbf{ride\_task}\\validity gate,\\Haversine, accumulate};
\node[tk] (t4) at (9.6,-0.2) {\textbf{crash\_task}\\impact $\rightarrow$ tilt\\$\rightarrow$ stillness};

% outputs
\node[snk] (t5) at (9.6,-2.5) {\textbf{ble\_task}\\GATT notify};
\node[snk] (nvs) at (4.6,-2.5) {\textbf{NVS buffer}\\deferred records};

\draw[ar] (gin) -- (t1);
\draw[ar] (iin) -- (t2);
\draw[ar,dashed] (ain) -- (t2);
\draw[ar] (t1) -- (q1);
\draw[ar] (t2) -- (q2);
\draw[ar] (q1) -- (t3);
\draw[ar] (q2) -- (t4);
\draw[ar] (t1.south) -- node[right,font=\sffamily\tiny,pos=0.55]{PPS} (t2.north);
\draw[ar] (t3) -- (t4);
\draw[ar] (t4) -- (t5);
\draw[ar] (t3.east) -- ++(0.7,0) |- (t5.east);
\draw[ar] ([yshift=2.2mm]t5.west) -- node[above,font=\sffamily\tiny]{no client}
          ([yshift=2.2mm]nvs.east);
\draw[ar] ([yshift=-2.2mm]nvs.east) -- node[below,font=\sffamily\tiny]{on connect}
          ([yshift=-2.2mm]t5.west);
\end{tikzpicture}
\caption{Data Processing Flow within the ESP32-S3 Microcontroller}
\end{figure}

\vspace{0.5em}
\noindent\textbf{(i) Data Input}

Two independent input streams enter the microcontroller. The GNSS receiver delivers NMEA
sentences as a byte stream on a hardware UART, typically at one update per second, together
with a one-pulse-per-second edge on a GPIO. The inertial measurement unit accumulates
samples in its FIFO at a substantially higher rate and is read in blocks over SPI. Where
fitted, the ADXL345 contributes only an interrupt edge, consuming no bus bandwidth until an
event occurs.

The disparity in rates is the reason the two streams are acquired by separate tasks. The
inertial stream must be drained promptly to avoid FIFO overrun, whereas GNSS sentences
arrive slowly but in bursts; handling both in a single loop would risk losing motion samples
precisely during the transient that matters most.

\vspace{0.5em}
\noindent\textbf{(ii) Core Processing}

\begin{enumerate}[leftmargin=*,itemsep=0.45em,topsep=0.5em]
\item \textbf{Sentence Parsing and Fix Validation.} The GNSS task reassembles sentences,
verifies checksums and extracts the navigation fields. Fixes reported invalid, or implying
an impossible speed relative to the previous accepted fix, are rejected before they can
affect the odometer.

\item \textbf{Sample Conditioning.} The inertial task converts raw counts to physical units,
subtracts the biases established during calibration and applies a short moving average. The
filter length is a deliberate compromise: long enough to suppress vibration-induced noise,
short enough to preserve the rise time of a genuine impact.

\item \textbf{Ride Computation.} The ride task evaluates the Haversine increment between
accepted fixes, accumulates distance and duration, tracks maximum and mean speed, and
appends coordinates to the route sequence at a decimated rate appropriate to map display.

\item \textbf{Crash Evaluation.} The crash task maintains a small state machine. From the
idle state, an impact above threshold moves it to a tilt-watch state; a tilt sustained
beyond the orientation bound moves it to a stillness-watch state; stillness maintained for
the required interval declares a crash. Failure of any condition within its window returns
the machine to idle, which is the mechanism by which potholes and normal cornering are
rejected.

\item \textbf{Record Formatting.} Results are packed into a compact binary record carrying
timestamp, coordinates, speed, cumulative distance, and status flags, sized to fit within a
single BLE notification so that no record requires fragmentation.
\end{enumerate}

\vspace{0.5em}
\noindent\textbf{(iii) Communication and Output}

\begin{enumerate}[leftmargin=*,itemsep=0.45em,topsep=0.5em]
\item \textbf{BLE Connection Management.} The BLE task advertises the custom service,
accepts a connection from the bonded handset, and manages subscription state for each
characteristic. Connection parameters are negotiated towards a longer interval once a ride
is under way, trading latency for power, and towards a short interval when a crash event is
pending.

\item \textbf{Publication and Deferred Delivery.} With a client connected, records are
notified as they are produced. With no client connected, they are appended to the
non-volatile buffer instead. On the next connection the buffer is drained in order before
live notification resumes, so the handset receives a continuous history.

\item \textbf{Error Handling and Reliability.} The firmware monitors for FIFO overrun, UART
framing errors, loss of satellite fix and BLE disconnection. Loss of fix suspends distance
accumulation but does not suspend crash detection, since the inertial path is independent of
positioning. Buffer writes are wear-levelled and record-framed, so that an unexpected power
loss can corrupt at most the record being written and never the records already committed.
\end{enumerate}

\subsection{Sequence Diagram}

The sequence proceeds in four phases: ride start, the logging loop, a crash event, and ride
end.

\begin{figure}[H]
\centering
\begin{tikzpicture}[
  font=\sffamily\tiny,
  ll/.style={draw,align=center,inner sep=2.5pt,minimum height=6mm,text width=15mm,
             fill=gray!10},
  ar/.style={-{Latex[length=1.6mm]}},
  x=1cm,y=1cm
]
% lifeline heads
\node[ll] (a) at (0,0)    {Rider};
\node[ll] (b) at (2.3,0)  {ATGM336H-5N};
\node[ll] (c) at (4.6,0)  {MPU6500};
\node[ll] (d) at (6.9,0)  {ESP32-S3};
\node[ll] (e) at (9.2,0)  {Flutter App};
\node[ll] (f) at (11.3,0) {SQLite};

\foreach \n in {a,b,c,d,e,f} { \draw[dashed,gray] (\n.south) -- ++(0,-9.3); }

\begin{scope}[every node/.append style={fill=white,inner sep=1.4pt}]
% phase labels
\node[anchor=west,font=\sffamily\tiny\bfseries] at (-0.85,-0.75) {Ride start};
\draw[ar] (0,-1.15)  -- node[above]{power on / start ride} (6.9,-1.15);
\draw[ar] (6.9,-1.5) -- node[above]{configure, read FIFO} (4.6,-1.5);
\draw[ar] (6.9,-1.85)-- node[above]{enable NMEA, 1\,PPS} (2.3,-1.85);
\draw[ar] (9.2,-2.2) -- node[above]{connect, subscribe (BLE)} (6.9,-2.2);

\node[anchor=west,font=\sffamily\tiny\bfseries] at (-0.85,-2.75) {Logging loop};
\draw[ar] (2.3,-3.1) -- node[above]{fix: lat, lon, speed} (6.9,-3.1);
\draw[ar] (4.6,-3.45)-- node[above]{$a_x,a_y,a_z,\omega$} (6.9,-3.45);
\draw[ar] (6.9,-3.8) -- node[above,align=center]{notify: speed, distance, route point} (9.2,-3.8);
\draw[ar] (9.2,-4.15)-- node[above]{append route point} (11.3,-4.15);
\draw[ar] (9.2,-4.5) -- node[above]{dashboard, map} (0,-4.5);
\node[anchor=west,font=\sffamily\tiny] at (7.05,-5.0)
      {\textit{loop: every GNSS update}};

\node[anchor=west,font=\sffamily\tiny\bfseries] at (-0.85,-5.55) {Crash event};
\draw[ar] (4.6,-5.9) -- node[above]{impact: $a_r > a_{\mathrm{th}}$} (6.9,-5.9);
\draw[ar] (4.6,-6.25)-- node[above]{sustained tilt $> \theta_{\mathrm{th}}$} (6.9,-6.25);
\draw[ar] (4.6,-6.6) -- node[above]{stillness confirmed} (6.9,-6.6);
\draw[ar] (6.9,-6.95)-- node[above]{notify: crash event} (9.2,-6.95);
\draw[ar] (9.2,-7.3) -- node[above]{alert + countdown} (0,-7.3);
\draw[ar] (0,-7.65)  -- node[above]{cancel (if false alarm)} (9.2,-7.65);
\draw[ar] (9.2,-8.0) -- node[above]{log crash / cancellation} (11.3,-8.0);

\node[anchor=west,font=\sffamily\tiny\bfseries] at (-0.85,-8.5) {Ride end};
\draw[ar] (6.9,-8.85)-- node[above]{final statistics} (9.2,-8.85);
\draw[ar] (9.2,-9.2) -- node[above,align=center]{commit ride, update odometer, reminder} (11.3,-9.2);
\end{scope}
\end{tikzpicture}
\caption{Sequence Diagram of the SmartRide Ride Logging and Crash Detection System}
\end{figure}


\vspace{0.5em}
\noindent\textbf{(i) Ride Start Phase.} On power-up the ESP32-S3 configures both sensors ---
setting the inertial measurement unit's ranges, filter bandwidth and FIFO, and enabling the
required NMEA sentences and the 1\,PPS output on the GNSS receiver --- and begins
advertising its BLE service. The handset, already bonded, connects and subscribes to the
ride and crash characteristics. The unit begins logging immediately, whether or not the
handset has connected.

\vspace{0.5em}
\noindent\textbf{(ii) Logging Loop Phase.} For each GNSS update the receiver delivers a fix,
which the microcontroller validates and converts into a distance increment and a route point.
The inertial measurement unit meanwhile delivers motion samples at a much higher rate. The
microcontroller notifies the current statistics to the application, which appends the route
point to SQLite and refreshes the dashboard and map for the rider. This loop repeats for the
duration of the ride.

\vspace{0.5em}
\noindent\textbf{(iii) Crash Event Phase.} The three conditions of the crash rule arrive in
sequence from the inertial stream: an impact above threshold, a tilt beyond the orientation
bound sustained for the required interval, and a period of stillness. Only on the third does
the microcontroller emit a crash event. The application raises an alert with a countdown
which the rider may cancel; the event and any cancellation are written to SQLite, both for
the rider's record and for subsequent threshold review.

\vspace{0.5em}
\noindent\textbf{(iv) Ride End Phase.} On detecting sustained stillness, or on explicit
instruction from the rider, the microcontroller emits the final ride statistics. The
application commits the ride, updates the cumulative odometer and, if the servicing interval
has been crossed, raises a maintenance reminder.

\section{WORK PLAN}

The preparation phase of the project runs from June to October 2026, covering topic
finalisation and literature review through the economic feasibility study and component
selection to the interim evaluation and the planning of the implementation phase. Each
column in the chart below represents half a month.

\begin{figure}[htbp]
\centering
\begin{tikzpicture}[font=\sffamily\scriptsize,x=0.85cm,y=0.64cm]
% month grid : 10 half-month columns, Jun..Oct
\foreach \i in {0,...,10} { \draw[gray!45] (\i,0) -- (\i,-8.5); }
\foreach \m/\lab in {1/Jun,3/Jul,5/Aug,7/Sep,9/Oct} {
  \node[font=\sffamily\small\bfseries] at (\m,0.45) {\lab};
}
\draw[gray!70] (0,0) -- (10,0);
\node[font=\sffamily\small\itshape] at (5,1.1) {Academic Year 2026};

% tasks : {row}/{start}/{end}/{label}
\foreach \r/\s/\e/\lab in {
  1/0/2/{Topic selection and finalisation},
  2/1/4/{Literature review},
  3/3/5/{Requirement analysis},
  4/4/6/{Economic feasibility study},
  5/5/7/{Component selection and costing},
  6/6/8/{System design and architecture},
  7/7/9/{Interim evaluation and documentation},
  8/8/10/{Planning for implementation phase}
} {
  \fill[blue!65!black] (\s,-\r+0.28) rectangle (\e,-\r-0.28);
  \node[anchor=east,font=\sffamily\scriptsize] at (-0.18,-\r) {\lab};
}
\end{tikzpicture}
\caption{Initial Work Plan of the SmartRide Project}
\end{figure}

The chart reflects the actual progression of the preparation phase. Topic selection and
literature review dominated the early months; requirement analysis and the economic
feasibility study followed; component selection and costing in August produced the revised
hardware set described in this report --- the ESP32-S3 in place of the original ESP32, the
ATGM336H-5N in place of the NEO-6M, the MPU6500 in place of the MPU6050, a managed
lithium-ion supply in place of a buck converter from the vehicle battery, and the ADXL345 as
an optional redundant impact channel. System design and architecture occupied September and
concluded with the interim evaluation, from which the implementation phase is now planned.

% ============================================================
\chapter{REFERENCES}

\begin{singlespace}
\begin{enumerate}[label={[\arabic*]},leftmargin=*,itemsep=0.85em,topsep=0.5em]

\item H. Alai, R. Rajamani, and K. Sanjaya, (2024), ``A Smart E-Scooter with Embedded
Estimation of Rear Vehicle Trajectories for Rider Protection,'' \textit{Mechanical Systems
and Signal Processing}, Elsevier.

\item G. Chhabra, et al., (2024), ``Internet of Things Based Smart Framework for the Safe
Driving Experience of Two Wheelers,'' \textit{Scientific Reports}, vol. 14, Nature
Publishing Group.

\item ``Development of Real-Time GPS Tracking and ETA on an Android Application Using ESP32
and u-blox NEO6MV2,'' (2024), in \textit{Proc. IEEE Int. Conf. on Computing, Communication
and Networking Technologies (ICCCNT)}, IEEE.

\item Espressif Systems, (2024), \textit{ESP32-S3 Technical Reference Manual}, Espressif
Systems (Shanghai) Co., Ltd.

\item Espressif Systems, (2024), \textit{ESP32-S3 Series Datasheet}, Espressif Systems
(Shanghai) Co., Ltd.

\item InvenSense Inc., (2014), \textit{MPU-6500 Product Specification}, Rev.~1.1, InvenSense
Inc., San Jose, CA.

\item Analog Devices Inc., (2015), \textit{ADXL345 Digital Accelerometer Data Sheet},
Rev.~E, Analog Devices Inc., Norwood, MA.

\item Zhongke Microelectronics, \textit{ATGM336H Series GNSS Module Data Sheet} (AT6558
multi-constellation positioning chip), Zhongke Micro (ZKW), Shanghai.

\item Bluetooth SIG, (2023), \textit{Bluetooth Core Specification, Version 5.4}, Bluetooth
Special Interest Group --- Generic Attribute Profile (GATT).

\item National Marine Electronics Association, \textit{NMEA 0183 Standard for Interfacing
Marine Electronic Devices}, Version 4.11, NMEA.

\item Google, (2024), \textit{Flutter Documentation}, \url{https://docs.flutter.dev}.

\item SQLite Consortium, (2024), \textit{SQLite Documentation},
\url{https://www.sqlite.org/docs.html}.

\item Ministry of Road Transport and Highways, Government of India, \textit{Road Accidents
in India}, Annual Report, MoRTH, New Delhi.

\item United Nations, (2015), \textit{Transforming Our World: The 2030 Agenda for
Sustainable Development}, United Nations General Assembly Resolution A/RES/70/1 ---
Targets 3.6, 11.2.

\item R. W. Sinnott, (1984), ``Virtues of the Haversine,'' \textit{Sky and Telescope},
vol.~68, no.~2, p.~159.

\end{enumerate}
\end{singlespace}

\end{document}
